\Article{32}
Wie die magnetischen Flüssigkeiten im Erdinneren tatsächlich verteilt sind, bleibt notwendigerweise unbestimmt. Nach einem allgemeinen Theorem, das bereits in Artikel 2 der \textit{Intensitas} erwähnt wird und bei einer anderen Gelegenheit ausführlich behandelt werden soll, lässt sich nämlich jede beliebige Verteilung der magnetischen Flüssigkeiten innerhalb eines dreidimensionalen Raumes stets durch eine Verteilung auf seiner Oberfläche ersetzen, sodass die Wirkung in jedem Punkt des Außenraums genau dieselbe bleibt. Daraus folgt leicht, dass \textit{ein und dieselbe} Wirkung im gesamten Außenraum aus unendlich vielen \textit{verschiedenen} Verteilungen der magnetischen Flüssigkeiten im Inneren hervorgehen kann.

Wir können dagegen die fiktive Verteilung auf der Erdoberfläche angeben, die der tatsächlichen Verteilung im Inneren hinsichtlich der daraus nach außen entstehenden Kräfte vollkommen gleichwertig ist. Wegen der \SourceLeaf{170}Kugelform der Erde können wir dies sogar auf sehr einfache Weise tun. Die Dichte des magnetischen Fluids in jedem Punkt der Erdoberfläche. das heißt die Menge des Fluids pro Flächeneinheit, wird nämlich durch die Formel
\[
\frac1{4\pi}(\frac VR-2Z)
\]
ausgedrückt oder durch
\[
-\frac1{4\pi}(3P'+5P''+7P'''+9P''''+\text{u. s. w.})
\]
Der Wert dieser Formel wird demnächst durch eine grafische Darstellung veranschaulicht. Hier sei nur bemerkt, dass er auf der nördlichen Hälfte der Erde negativ und auf der südlichen positiv ist, wobei die Trennlinie den Äquator zweimal schneidet (bei $6^0$ und $186^0$ Länge) und sich auf beiden Seiten bis zu etwa $15^0$ nördlicher beziehungsweise südlicher Breite von ihm entfernt. Ferner gibt es auf der nördlichen Hälfte \textit{zwei} Minima und auf der südlichen nur ein Maximum. Nach einer überschlägigen Rechnung ergeben sich für diese Minima und das Maximum
\begin{center}
\begin{tabular}{rlrlrl}
$-209{,}1$&in&$55^0$&N. Breite&$263^0$&Länge\\
$-200{,}0$&in&$71^0$&N. Breite&$116^0$&Länge\\
$+277{,}7$&in&$70^0$&S. Breite&$154^0$&Länge
\end{tabular}
\end{center}
Den Werten selbst liegt die Einheit unserer Elemente zugrunde. Sollen sie in absolutem Maß ausgedrückt werden, müssen sie daher noch mit $0{,}0034941$ multipliziert werden.

\Article{33}
Unsere Elemente sollen, wie oben bereits gesagt wurde, nur als erste Näherung gelten. Als solche stimmen sie nach Artikel 29 hinreichend gut mit den Beobachtungen überein. Es besteht kein Zweifel, dass eine auf diesen Beobachtungen beruhende Verbesserungsrechnung eine viel größere Übereinstimmung erzielen würde. Die einzige Schwierigkeit einer solchen Rechnung wäre ihr Umfang. Dieser bliebe immer noch abschreckend groß, selbst wenn man zur Abkürzung ähnliche Verfahren anwenden wollte, wie sie die Astronomen zur Verbesserung der Bahnelemente von Planeten und Kometen verwenden. Diese Schwierigkeit ließe sich zwar leicht
\SourceLeaf{171}
überwinden, wenn die Arbeit auf mehrere Personen verteilt werden könnte, die die Rechnungen ausführen. Dennoch dürfte es kaum ratsam sein, eine solche Verbesserung schon jetzt vorzunehmen, solange die Daten von so vielen Orten, die man unbedingt einbeziehen müsste, noch so unzuverlässig sind. Am besten wird man vorerst die Elemente weiter mit Beobachtungen vergleichen. So lässt sich ein Weg finden, den allgemeinen Karten eine viel größere Zuverlässigkeit zu geben, als es mit dem bisher ausschließlich empirischen Verfahren möglich war. Es sei uns aber gestattet, einige Blicke auf die künftigen Fortschritte der Theorie zu werfen, deren vollständige Verwirklichung freilich noch in weiter Ferne liegen mag.

\Article{34}
Um die Elemente zufriedenstellend zu verfeinern und zu vervollständigen, müssen die Beobachtungsdaten viel höhere Anforderungen erfüllen als bisher. Sie sollten an allen herangezogenen Punkten eine Genauigkeit aufweisen, die bislang nur an ganz wenigen erreicht wurde. Sie sollten von den unregelmäßigen Bewegungen bereinigt sein und für ein und denselben Zeitpunkt gelten. Es wird noch lange dauern, bis solche Anforderungen erfüllt werden können. Am dringendsten sind zunächst \textit{vollständige} (das heißt alle drei Elemente umfassende) Beobachtungen an dem einen oder anderen Punkt in jenen großen Gebieten, in denen derartige Beobachtungen bislang ganz fehlen. Tatsächlich ist jeder neu hinzukommende Punkt für die allgemeine Theorie umso wichtiger, je weiter er von den anderen Punkten entfernt liegt, über die wir bereits verfügen.

Nach einer hinreichenden Zwischenzeit wird man die Elemente für einen zweiten Zeitpunkt erneut bestimmen und so ihre säkularen Änderungen erhalten. Dazu wird es aber offensichtlich unerlässlich sein, das bisherige Maß für die Intensitäten ganz aufzugeben und durch ein absolutes Maß zu ersetzen.

Im Laufe der kommenden Jahrhunderte werden auch diese Änderungen nicht mehr gleichförmig erscheinen, und die Erforschung des Verlaufs der Elemente wird den Naturforschern unerschöpflichen Stoff für Untersuchungen bieten.

\SourceLeaf{172}
\Article{35}
Die Zukunft wird aber auch Aufschluss über interessante Punkte der Theorie geben.

In unserer Theorie wird angenommen, dass jeder messbare magnetisierte Teil des Erdkörpers genau so viel positives wie negatives Fluid enthält. Hätten die magnetischen Flüssigkeiten keinerlei Wirklichkeit und wären sie nur ein fiktiver Ersatz für galvanische Ströme in den kleinsten Teilen der Erde, so ergäbe sich diese Gleichheit schon daraus, dass diese Ersetzung zulässig ist. Schreibt man den magnetischen Flüssigkeiten hingegen eine wirkliche Existenz zu, könnte man, ohne sich zu widersprechen, die vollkommene Gleichheit der Mengen beider Flüssigkeiten bezweifeln. Bei einzelnen magnetischen Körpern (natürlichen oder künstlichen Magneten) ließe sich durch sehr genaue Versuche leicht entscheiden, ob sie einen merklichen Überschuss der einen oder der anderen Flüssigkeit enthalten. Im ersten Fall müsste ein mit einem solchen Körper belasteter Lotfaden von der Senkrechten abweichen (und zwar in Richtung des magnetischen Meridians). Würden solche Versuche mit vielen künstlichen Magneten an einem hinreichend weit von Eisen entfernten Ort niemals die geringste Abweichung zeigen (was wohl zu erwarten ist), so wäre diese Gleichheit zwar auch für die ganze Erde mit größter Wahrscheinlichkeit anzunehmen, doch bliebe die Möglichkeit einer gewissen Ungleichheit bestehen.

Eine solche Ungleichheit würde in unserer Theorie nur den Unterschied bewirken, dass $P^0$ (Artikel 17) nicht mehr $=0$ sein würde Die Folge wäre, dass im gesamten unendlichen Außenraum dem Ausdruck für $Z$ noch das Glied $\frac{RRP^0}{rr}$ und damit auf der Erdoberfläche das (konstante) Glied $P^0$ hinzugefügt werden müsste, während $X$ und $Y$ davon völlig unberührt blieben. Wenn künftig sehr viel mehr genaue Beobachtungen vorliegen als heute, wird sich tatsächlich ermitteln lassen, ob ihre genaue Darstellung einen von null verschiedenen Wert für $P^0$ erfordert oder nicht. Beim gegenwärtigen Stand der Daten könnte ein solches Unternehmen aber noch gar kein Ergebnis liefern.
